% ================================================================
% DATA AND DESCRIPTIVE EVIDENCE
% ================================================================

% ==================================================================
% 3.1 INSTITUTIONAL CONTEXT
% ==================================================================
\subsection{The EU Emissions Trading System}\label{subsec:euets}

% TODO: Describe the EU ETS:
%   - European-wide cap-and-trade for GHG emissions, launched 2005
%   - Sectoral targeting: ~60 NACE 4-digit categories (out of 732)
%     covering energy + heavy manufacturing (cement, steel, etc.)
%   - Covers ~45% of EU emissions, ~5% of global emissions
%   - Within-sector firm-level targeting via participation thresholds
%     (e.g., installations above certain capacity/output thresholds)
%
% Note: our sufficient statistic is derived for a carbon tax, but the
% EU ETS can be interpreted as one for the purpose of first-order
% comparative statics (given a price).


% ==================================================================
% 3.2 EU ETS COVERAGE IN BELGIUM
% ==================================================================
\subsection{EU ETS Coverage in Belgium}\label{subsec:coverage}

% TODO: Show how much of Belgian emissions the EU ETS covers:
%   - Aggregate: ~60% of emissions from stationary sources, ~40% of
%     total Belgian emissions
%   - By sector: present a table or figure showing coverage rates
%     across NACE sectors in Belgium
%   - Reference the sample coverage table (currently in slides appendix)
%     showing count of firms, GDP coverage, wage bill coverage, and
%     emissions coverage for selected years (2002, 2012, 2022)


% ==================================================================
% 3.3 DATA SOURCES
% ==================================================================
\subsection{Data}\label{subsec:data_sources}

% TODO: Describe each data source:
%
% 1. EUTL (EU Transaction Log):
%    - Installation-year-level data on verified emissions, firm ownership,
%      4-digit sector classification
%    - Covers all installations targeted by EU ETS, 2005-2022
%    - We aggregate to firm-year level using ownership information
%
% 2. NBB firm-level micro data (2002-2022):
%    - Business-to-business (B2B) transactions: firm-to-firm trade flows
%      => construct firm-year-level I-O matrix Omega for the Belgian
%      private sector
%    - Annual accounts (balance sheets): firm-level value added, wage
%      bill, employment, assets
%
% 3. Energy balances and national inventories:
%    - Sector-level emissions data used for imputation of non-targeted
%      firms' emissions
%
% Summary: we directly observe tau (targeting), E (emissions for targeted
% firms), and Omega (I-O matrix) => can compute Psi and Psi_e.


% ==================================================================
% 3.4 EMISSIONS IMPUTATION
% ==================================================================
\subsection{Imputation of Non-Targeted Firms' Emissions}\label{subsec:imputation}

% TODO: For firms not in the EU ETS, we impute emissions:
%
% Step 1: Obtain sector-level total emissions Z_{s,t} from energy
% balances + national inventories.
%
% Step 2: For each sector s, compute non-targeted emissions as
% Z_{s,t} - sum_{j in EUETS} z_{j,s,t}
%
% Step 3: Distribute non-targeted emissions proportional to firm size:
%   z_{i(s),t} = [x_{i(s),t} / sum_{j(s) not in EUETS} x_{j,s,t}]
%                * (Z_{s,t} - sum_{j(s) in EUETS} z_{j,s,t})
%
% Validation: show observed vs imputed emissions figure for firms where
% we have both (R^2 = 0.56). Discuss plans to improve using firm-level
% input bundle and product offering data.
%
% Include figure: output of imputation validation
% (macro_lunch_may25/figures/obs_vs_imputed_emissions.png)


% ==================================================================
% 3.5 DESCRIPTIVE FACTS
% ==================================================================
\subsection{Descriptive Facts}\label{subsec:facts}

% ------------------------------------------------------------------
% 3.5.1 Emissions decomposition in Belgium
% ------------------------------------------------------------------
\subsubsection{Statistical Decomposition of Belgian Emissions}

% TODO: Apply the scale/technique/composition decomposition to the
% observed fall in Belgian emissions between [start year] and [end year].
%
% This is a purely statistical exercise (no model needed): using the
% observed data on emissions, output, and emission intensities, decompose
% the aggregate emissions change.
%
% Present results in a table or figure showing the contribution of each
% channel over time.

% ------------------------------------------------------------------
% 3.5.2 Within-sector dispersion in emission intensity
% ------------------------------------------------------------------
\subsubsection{Dispersion in Emission Intensity}

% TODO: Fact 1 — emission intensity is vastly heterogeneous across firms
% within narrowly-defined sectors.
%
% For each firm i in year t, define:
%   emissions productivity := log(1/e_{i,t})
%   labor productivity     := log(1/Omega_{i,l,t})
%   capital productivity   := log(1/Omega_{i,k,t})
%
% Measure dispersion as 90th - 10th percentile within industry.
% Take average across industries.
%
% Show figure comparing dispersion in emissions productivity vs labor
% and capital productivity (adapt from slides: avg_dispersion.png).
%
% Show separately for EUETS firms only and for all Belgian firms.

% ------------------------------------------------------------------
% 3.5.3 Within-sector dispersion in network-adjusted emission intensity
% ------------------------------------------------------------------
\subsubsection{Dispersion in Network-Adjusted Emission Intensity}

% TODO: Fact 2 — network-adjusted emission intensity (Psi_e) is also
% heterogeneous across firms within sector.
%
% This is a novel fact: even if two firms in the same sector have similar
% scope 1 emissions, their network-adjusted emission intensity can differ
% substantially because they source inputs from supply chains with
% different emission profiles.
%
% Present dispersion statistics and/or figures analogous to Fact 1.
